\section{Cuts, comparison, and spatial estimates}
\label{sec:spatial}

In \crefrange{sec:spatial}{sec:gaussian}, heights are measured in the
normalization $K_\delta=2H_\delta$, whose nearest corner increments have
magnitude one.

The field limit needs control of local patterns, whereas the cluster
limit also needs estimates near the irregular boundaries left by an
exploration. We obtain both from one edge representation. Fix one of the
two currents and call it $N$. Let $W$ be its exact edge complement on the
other lattice: a $W$ edge is open precisely when the crossing $N$ edge is
closed. Write $s$ for the spin on the $W$ lattice and $d$ for the spin on
the $N$ lattice. Thus $W$ is monochromatic in $s$, while $N$ is
monochromatic in $d$. The graph $W$ is generally different from the other
current in the coupling. At an ambiguous edge its opening probability is
$t$, not $1-t$.

A \emph{flat chamber} means a component cut off by an open monochromatic
$W$ perimeter, with its height on that perimeter fixed. Boundary walks
may repeat vertices. At a pinch, the two incident faces on the unsearched
side remain separate unless they are actually connected there. One can
equivalently use medial boundary polygons displaced by one lattice unit
into the chamber. This convention changes Euclidean distances by
$O(\delta)$ and never joins two chambers across a pinch.

We use three kinds of conditional information. An \emph{exact cut
history} may contain the chosen current $N$, its colors, and the heights
on the revealed side of completed $W$ perimeters. The edge factor below
removes that information from every residual chamber. A comparison or
annular search \emph{inside} a residual chamber reveals only $(s,W)$,
and reveals an edge only after its endpoint spins. Finally, an auxiliary
matching fill changes the removed side of a cut for a geometric circuit
test; it is used through its accessible marginal and is not extra
conditioning of the original law. The edge decoration of the other
current has been summed out. It is never part of the stopping data on a
cut or on its unsearched side: conditioning it on an open $W$ edge can
restore an interaction between the transverse spins. Data strictly on an
already separated side may be retained by the exact cut factorization.

All constants below may depend on the fixed parameter $t\in[1/2,1)$ and
on stated macroscopic buffer ratios. They do not depend on the mesh,
on the number of microscopic edges of a prescribed pattern, or on the
length or number of fingers of a retained perimeter. In comparisons
involving several independently retained source wires, dependence on
their number will always be stated.

\subsection{Exact cuts and the comparison order}

\begin{lemma}[Finite cuts and colors]\label{lem:finite-cuts}
For every finite domain with either of the model's flat boundary
conditions, the joint weight of $(s,W,d)$ is a product
of the following factors, one for each mutually dual interior edge pair:
\begin{equation}\label{eq:spatial-edge-weight}
 \begin{cases}
 t\ind_{\{s_i=s_j\}},& W\text{ open},\\
 \bigl(t\ind_{\{s_i\ne s_j\}}+
       (1-t)\ind_{\{s_i=s_j\}}\bigr)
       \ind_{\{d_i=d_j\}},& W\text{ closed}.
 \end{cases}
\end{equation}
The subscripts in each indicator refer to the endpoints on its own
lattice. In particular:
\begin{enumerate}
\item Conditional on $(s,W)$, the free $N$ components receive independent
fair $d$ signs. A prescribed wired $N$ component retains its prescribed
sign.
\item An open monochromatic $W$ cut fixes a constant height on its
perimeter and carries no transverse-spin interaction. Conditional on
the cut and on an exact cut history on its other side, different
unsearched chambers have independent flat laws, with free $N$ colors.
\item On the $s$ lattice, height differences from an exterior wall are
sums of jumps of size $2$. Each free $N$ cluster contributes one fair
relative sign, common to all its odd holes, and contributes zero across
its even holes. These signs are independent conditional on $(N,s)$.
Here a hole is odd when its adjacent $s$ spin differs from the adjacent
exterior $s$ spin of that cluster.
\item These statements hold at stopping explorations which reveal whole
$N$ clusters met by a connected growing seed attached to the exterior,
or stop after a prescribed number of successive odd rims. In the wired
$N$ case the first exposed object is its boundary cluster; its holes
have opposite-parity $W$ rims at heights $+1$ or $-1$ relative to the
initial height, and their same-lattice currents are free.
\end{enumerate}
\end{lemma}

\begin{proof}
If $s$ disagrees, the crossing $N$ edge is forced open and its $d$ spins
agree. If $s$ agrees and $W$ is closed, the weight is $1-t$ and the same
agreement is required. Finally, if $W$ is open, summing the possibilities
for the other current gives weight $t$ whether the transverse spins agree
or disagree. This proves \eqref{eq:spatial-edge-weight}. Its dependence on
$d$ consists exactly of equality constraints on $N$ edges, proving the
first assertion. Fixing a wired component's sign removes the constant
factor $2$ associated with that component, independent of its geometry.

On an open $W$ perimeter only the first line of
\eqref{eq:spatial-edge-weight} occurs. The transverse variables on its
two sides have no interaction across it. The remaining product factors
split by chamber. Along the perimeter $s$ is constant, and the corner
height rule gives constant height between successive same-parity
vertices. This proves the second assertion, first for a simple perimeter
and then for a boundary walk by keeping its incident faces separate.

For two $s$ vertices $v,v'$ incident to a transverse vertex $f$,
\begin{equation}\label{eq:spatial-height-jump}
 K(v')-K(v)=\varepsilon_W d(f)\bigl(s(v)-s(v')\bigr),
\end{equation}
where $\varepsilon_W=1$ when the $s$ lattice is primal and
$\varepsilon_W=-1$ when it is dual. This fixed orientation sign has
no effect on the independent fair relative signs.
The rim facing any specified hole of a connected $N$ component belongs
to one incident $W$ component, as does its exterior rim. Thus crossing
that $N$ component gives either zero or $2$ times its $d$ color times
the fixed exterior $s$ sign. All odd holes give the same relative jump.
The third assertion now follows from the first.

For the stopping assertion, inspect no variable beyond a completed open
$W$ rim sealing an unsearched face. Whenever the search queries $(s,W)$,
it exposes endpoints before edge states. For each possible finite
history, its event is the cylinder of the values queried by the
deterministic search. The product factors inside the residual chambers
are precisely those already described, even if the rule deciding which
completed cluster to examine next used its revealed $d$ color. Thus the
conditional residual kernels, and not just their average, are the stated
independent flat laws.
Skipping an even hole reveals no new jump; stopping at an odd rim
retains its observed jump. A connected exterior-attached seed and all
clusters it has met form one connected explored region $S$. Every edge
from $S$ to an unsearched vertex is $N$-closed, since every met cluster
has been revealed in full. For one geometric component $C$ of the
unsearched vertices, its complement is connected through $S$; every
other unsearched component attaches to $S$ in the underlying graph.
Thus the edge boundary of $C$ is a planar bond and its dual is one
open $W$ perimeter, with the stated medial resolution at pinches.
It is monochromatic, and all its remote connections are retained.
This argument concerns connectedness in the underlying graph;
the added seed edges need not themselves be $N$-open.
The seed may have visited several height levels, but this causes
no varying boundary value in a residual chamber: every edge on its
component-side $W$ perimeter joins equal $s$ spins, and
\eqref{eq:spatial-height-jump} gives zero height increment along that
edge. The connected perimeter therefore has one constant height.
The history determines that constant using only already revealed
clusters and their colors; no interior or future information is used.
For a wired current,
the incident opposite-parity height differs from the boundary height by
one corner increment. After its boundary cluster is removed, no interior
$N$ cluster is wired to the original boundary. This proves the last
assertion.
\end{proof}

The same cut argument permits a single exploration to interleave several
residual chambers behind an exposed wired layer. Attach each new seed
vertex to the common exposed region through already visited vertices or
known seed edges, with no unsearched free vertex in the attachment. The
union remains connected, and the unsearched components retain their
product of flat kernels. For example, a fixed global column sweep can
visit every earlier column in every chamber before the next one; a
missing neighbor is then in the exposed layer or is already visited.
This is the composite exploration used later.

The useful order is the coordinatewise order on
\begin{equation}\label{eq:spatial-order}
       (s,W^+,-W^-),
\end{equation}
where $W^+$ and $W^-$ are the open edges of the two spin colors, and
$-1<+1$ for vertex spins. An edge prescription is made only after its
endpoint spins have been prescribed. A retained plus wire in a lower
boundary condition must remain connected in the upper condition;
for minus wires the inclusion is reversed. We call such boundary
conditions \emph{ordered}. All other already specified connections
remain part of the conditional state, even if their paths run outside
the window being examined.

\begin{proposition}[Comparison with pins and wires]
\label{prop:comparison}
The finite $(s,W)$ law is positively associated in
\eqref{eq:spatial-order}, as is every conditional law obtained by fixing
vertex spins and compatible edge states with both endpoints fixed.
Increasing additional vertex and compatible edge prescriptions, or
changing retained connections according to the ordered boundary rule,
increases the residual law. The analogous statement holds with the two
colors reversed.

These comparisons permit monotone circuit couplings: search from the
changed data toward the region of interest, and stop at an open plus
circuit in the lower configuration, or an open minus circuit in the
upper configuration. Beyond such a common cut, the two configurations,
including complementary cluster colors, can be sampled identically.
\end{proposition}

\begin{proof}
Summing $d$ in \eqref{eq:spatial-edge-weight} contributes $2^{k(N)}$,
up to the constant just noted for a fixed wired color. Planar Euler
counting gives
\begin{equation}\label{eq:spatial-euler}
       2^{k(N)}=C\,2^{k(W)+|W|},
\end{equation}
where $C$ depends only on the underlying finite graph and its boundary
convention. Indeed the components of the complementary dual graph are
the faces of $W$, whose number is $|W|-|V(W)|+k(W)+1$.
For a local window, contract the actual same-color boundary contacts
joined by exterior $W$ paths and write $k_\xi(W)$ for the resulting
component count. The same identity holds with $k_\xi(W)$, up to a
constant depending on this boundary convention. These identifications
are retained in the law even when their paths lie far outside the window.

Put $r=(1-t)/t>0$. After division by the constant $t^{|E|}$, an open
monochromatic $W$ edge has weight $2$, a closed monochromatic edge has
weight $r$, and a disagreeing edge is closed with weight $1$.
Represent $2^{k(W)}$ by auxiliary colors $\rho\in\{-1,1\}$.
Summing the $W$ state on a monochromatic edge gives
\begin{equation}\label{eq:spatial-ferro-factor}
 r+2\ind_{\{\rho_i=\rho_j\}}
       =r+1+\rho_i\rho_j.
\end{equation}
Consequently the $s$ marginal is the product of a positive factor
$(1+r)^{|E_{\rm mono}|}$ and the even-subgraph partition functions on
the two monochromatic induced graphs, with edge parameter
$q=(1+r)^{-1}\in(0,1)$.

For completeness, let $Z_F(A)$ denote such an even-subgraph partition
function on an available edge set $A$, with ferromagnetic
identifications $F$ retained. It is increasing and log-supermodular in
$A$. One way to see both facts is to write it as
the zero-field ferromagnetic Ising partition function divided by its
single-edge $\cosh$ factors. Adding an edge multiplies the normalized
partition function by $1+q\E[\rho_i\rho_j]$. The expectation is
nonnegative and increases when other nonnegative couplings are added,
by the two Griffiths inequalities \cite{Griffiths1967I}; a complete
proof for nonnegative couplings is given in
\cite[Section~3.8.1, Theorem~3.49]{FriedliVelenik2017}.
Indeed differentiating an expectation with respect to a coupling gives
the corresponding nonnegative covariance. Identifications are limits
of positive bonds: after dividing by the common maximal exponential,
the finite-state factor is $e^{-2K D_F}$, where $D_F$ counts violated
identifications. As $K\to\infty$ it tends to their joint indicator.
Thus the same inequalities survive the finite-state limit. The
single-edge factors cancel from the mixed inequality, proving
\[
       Z_F(A\cup B)Z_F(A\cap B)\ge Z_F(A)Z_F(B).
\]
The same incremental-edge formula gives the boundary comparison that
will be needed below: if $A\subseteq B$ and $F\subseteq F'$, then
\begin{equation}\label{eq:spatial-mixed-gks}
       Z_{F'}(B)Z_F(A)\ge Z_{F'}(A)Z_F(B).
\end{equation}
Indeed each multiplier $1+q\E[\rho_i\rho_j]$ increases when the bonds
implementing $F'$ are added. Integrating the Griffiths covariance
inequality and then taking their infinite-bond limits proves
\eqref{eq:spatial-mixed-gks}.
For plus vertex sets $S,T$,
$E(S\cap T)=E(S)\cap E(T)$ and
$E(S\cup T)\supseteq E(S)\cup E(T)$.
Apply the preceding inequality and monotonicity to both the plus and
minus induced graphs. The factor $(1+r)^{|E_{\rm mono}|}$ is itself
ferromagnetic. This proves the spin lattice condition, to which the
finite association and comparison theorems
\cite{FortuinKasteleynGinibre1971,Holley1974} apply.

Conditional on $s$, the open edges on each monochromatic graph have
the ordinary random-cluster law with cluster weight $2$ and edge
weight $v=2/r$. Given the other edges, its opening probability is
$v/(v+2)$ if the endpoints are disconnected and $v/(v+1)$ if they
are already connected. The probabilities increase with the available
connections, so common-uniform edge updates give the edge comparison.
Together with the spin lattice
condition, this gives positive association and the order
\eqref{eq:spatial-order}. Fixed closed edges with prescribed endpoints
give constants; fixed open edges give ferromagnetic identifications.
Thus the same argument applies after all the allowed prescriptions.
The ordered-wiring comparison follows from
\eqref{eq:spatial-mixed-gks} for plus identifications and its reflected
version for minus identifications. In particular,
arbitrarily shaped exterior connections cannot be silently discarded.

Use these conditional comparisons at each step of a simultaneous
exploration, always preserving each configuration's residual
conditional law. A plus circuit in the lower configuration is also
open plus in the upper. By \cref{lem:finite-cuts}, the laws on the
sealed side then have the same constant boundary spin and no
transverse interaction across the circuit. Couple those laws
identically. No $N$ component crosses the cut, so its free colors may
also be coupled identically. The minus case is the reflected argument.
\end{proof}

There is also a useful bounded comparison when monotonicity is not the
issue. Suppose two local conditional weights have exactly the same
available edges, pins, and local edge factors, and differ only in the
partition of $k$ prescribed same-color source labels. Each label may
have many boundary contacts, which are kept identified in both laws.
For every configuration the two component counts differ by at most
$k-1$. After normalization, \eqref{eq:spatial-euler} therefore gives
\begin{equation}\label{eq:spatial-wire-density}
  2^{-2(k-1)}\le \frac{\dd\mu_\xi}{\dd\mu_{\xi'}}
                     \le 2^{2(k-1)} \qquad (k\ge1).
\end{equation}
This statement does not permit a single source to be split into its
many visible entrances, nor a change of local edge factors. A wire in
this calculation is an identification in a probability law, not a
geometric edge along which a circuit may travel.

At fixed mesh these conclusions also apply to bounded searches in the
plane. Apply them first in finite flat exhaustions and take a joint
subsequential limit of the configurations and of the finite search
transcripts. The limit is the common plane law obtained from flat
exhaustions in \cite[Theorem~17.1]{DKLM2026}. Equality on the sealed side
of a common
cut holds on every finite subset and hence on the whole countable
lattice in the limit, including the complementary colors. That theorem
also gives no unbounded complementary component and the fair color law
there. A failed circuit search stops at its prescribed buffer boundary;
it never conditions the residual law on a future success.

\subsection{Buffered circuits and rough walls}

We specify the geometric completion used when a test window meets a
removed region. A face of the $W$ lattice is centered at a vertex of
the dual $N$ lattice. For a revealed set of $N$ vertices and seed edges,
take the union of these closed faces and the faces along the seed. We
call this its \emph{source cell region}. It is within $O(\delta)$
of the drawn revealed graph. Holes occupied by unsearched chambers are
retained; this operation does not fill the hull of a cluster. Its
resolved complementary chambers and the medial perimeters above differ
by $O(\delta)$, including at pinches.

Let $F$ be such a removed cell region and $U$ an accessible component.
Suppose every frontier edge between $U$ and $F$ is a prescribed open
$W$ edge of one color $\sigma$. Define the \emph{augmented matching
graph} in a test window to consist of the actual $W^\sigma$ edges on the
accessible side and artificial matching $W$ edges contained in $F$.
The artificial edges are the lattice one-skeleton of the cell region;
the frontier edges themselves retain their actual prescribed states.
A circuit in this graph is an ordinary geometric circuit in the
window. It may use the artificial edges inside $F$. An underlying dual
path disjoint from $F$ cannot cross it without crossing an actual
$W^\sigma$ edge. In particular, an abstract remote wire is never used
as a geometric edge of this circuit.

There are two exact uses of this construction. For one component $U$,
complete each removed source on its other side by matching spins and
open edges. Keep the original conditions on any other boundary of $U$.
The frontier factors in \eqref{eq:spatial-edge-weight} are constants
$t$ with no transverse interaction. Integrating the removed variables
therefore shows that the $U$ marginal of this completion is exactly its
original conditional marginal. This remains true when $U$ has several
boundary components, provided all the completed frontiers have the
stated color and the test windows avoid any other prescribed pieces.

For the second use, suppose the residual chambers are the disjoint
complementary components of one common connected source cell region
$F$ attached to the exterior. Conditional on an exact cut history they
have a product of flat laws. Reflect $s$ in each chamber, if needed, so
that all its frontier spins have one color $\sigma$, and put matching
spins and open edges in $F$. The same factorization realizes exactly
this product as the accessible marginal of a single completed graph:
every factor on the removed side cancels on normalization. There is no
factor
depending on the number of chambers. We call this the
\emph{product class with a common fill}. It includes the chambers behind
an exposed wired layer
and subsequent composite seed explorations. An arbitrary collection of
independent chamber copies is not asserted to have such a realization.
Finally, when a completion changes only a partition of $k$ indivisible
source labels while preserving all local factors and pins, the comparison
\eqref{eq:spatial-wire-density} gives the stated bounded substitute for
exact equality. These are the only completion rules used below.

\begin{proposition}[Buffered circuits and attached-wall bounds]
\label{prop:circuits}
Fix positive aspect and buffer ratios. There are $p>0$ and a fixed
lattice cutoff $c_*>0$ with the following properties, uniformly for
every physical test radius $r\ge c_*\delta$ and every mesh. Test
neighborhoods must lie in the ambient domain, apart from matching
cell regions to which one of the preceding completion rules applies.
\begin{enumerate}
\item In a buffer containing no pins, a circuit of either prescribed
$W$ color in a fixed-aspect annulus has probability at least $p$,
conditional on arbitrary admissible $(s,W)$ data outside the buffer.
Compatible pins of the requested color inside the buffer only improve
the bound. Finite chains of overlapping circuit tests give a positive
lower bound for a fixed-aspect tube connection.
\item After any admissible $(s,W)$ stopping exploration, failure of
matching circuits on $m$ separated available scales has probability
at most $(1-p)^m$, when each scale is chosen from the current history
before its interior is examined. In particular, suppose a prescribed
connected
matching wall reaches from $B(z,r)$ to outside $B(z,R)$, with a fixed
buffer around $B(z,R)$. A path avoiding that wall and all attached
matching circuits, and joining $B(z,r)$ to outside $B(z,R)$, has
probability at most
\begin{equation}\label{eq:spatial-attached-arm}
       C\left(\frac{r+\delta}{R}\right)^\alpha
\end{equation}
for some $C,\alpha>0$.
\item When a test buffer meets $F$, the circuit and tube events mean
circuits and connections in the augmented matching graph just defined.
They have the same lower bounds under the exact completion of one
component or the exact completion with a common fill. A circuit
attached to $F$ blocks an accessible dual path by its actual part
together with the forbidden region $F$.
If only the $k$-label comparison is available, the probability bounds
lose at most $2^{O(k)}$ in total. The number of local entrances of one
label does not enter this cost.
\end{enumerate}
The bulk assertion in item~1 also holds before removing the model's
single wired $N$ component. Conditioning in a search concerns $(s,W)$;
an exact cut history with exterior colors is allowed because its
accessible marginal is already a fresh flat kernel.
\end{proposition}

\begin{proof}
The imported input is \cite[Theorem~16.4]{DKLM2026}: for
$\mathbf c=1/t\in[1,2]$, an opposite-color spin-percolation circuit in
a lattice annulus of radius at least $4$ has a positive lower bound,
uniform in its radius, under the worst constant wired boundary
condition. The decoration in that theorem opens an ambiguous edge
with probability $1/\mathbf c=t$, so it is precisely $W$. Its
exterior-conditioning version, \cite[Corollary~16.6]{DKLM2026},
allows arbitrary data of the same spin parity and its edges outside
the buffered ball. We only use this statement for one parity. Its
conditional FKG hypotheses agree with \cref{prop:comparison}.

After the fixed change from source lattice units to physical units,
the lower bound holds for $r\ge c_*\delta$. Comparison with the worst
constant boundary condition gives item~1 for squares; a fixed finite
cover gives the other stated aspect ratios and the tube version.
Intersect its matching circuit events by positive association.
Matching pins are favorable in \eqref{eq:spatial-order}. In the
wired $N$ law before removal, fixing the single wired $d$ color removed
the geometry-independent factor $2$ in the proof of
\cref{prop:comparison}. Its $(s,W)$ comparison is therefore the same.
For the unpinned bulk test, choose a deterministic circuit in the
unused part of the buffer between the annulus and the buffer boundary.
Compare to the law with that circuit pinned to the worst $W$ color.
Its sealed side is a $W$-flat kernel by \cref{lem:finite-cuts}, so
the same lower bound applies there despite the arbitrary data beyond
the buffer. Requested-color pins are then handled by monotonicity.
This proves the claimed wired bulk case.

For separated scales, reveal earlier trial neighborhoods before
testing the next one. The conditional failure probability remains
at most $1-p$, so multiplying conditional bounds proves item~2
without an independence assumption. Take a geometric sequence of
scales between $r+\delta$ and $R$, separated enough to have disjoint
buffers. A circuit at one such scale meets the wall because the wall
crosses the annulus. It separates the inner from the outer scale.
This gives \eqref{eq:spatial-attached-arm}; the $\delta$ term absorbs
the finitely many scales below the lattice cutoff.

For item~3, apply item~1 in the completed graph. All additional pins
in a test buffer have the requested matching color. The exact
accessible marginal or \eqref{eq:spatial-wire-density} transfers the
event to the original law. In this transfer an edge of the circuit
inside $F$ is explicitly artificial; every other edge of it is an
actual matching edge. Thus a full circuit is available even when an
accessible side is slit and has no full circuit by itself. The
geometric blocking assertion follows from the definition of the
augmented graph. For a connected source the boundary incident to one
fixed complementary component is one boundary walk. It can enter the
window arbitrarily often, but these entrances form one retained label.
This verifies exactly the premise on labels when that bounded
comparison is used.
\end{proof}

A useful consequence concerns a stopping set selected using revealed
colors. Under the preceding flat kernel, if a cluster meeting a seed
at distance at least $R$ from $B(z,r)$ reaches that ball, there is an
$N$ arm from $B(z,r)$ to outside $B(z,R)$. Testing the $W$ annuli in
that preceding $(s,W)$ marginal bounds this event by
\eqref{eq:spatial-attached-arm}, regardless of the later rule selecting
the cluster. Similarly, before revealing discovered colors, reaching a
ball before a fixed number of odd rims entails failure to find that
many alternating pairs of $W$ circuits on the intervening scales.
Conditional iteration gives a bound of the same polynomial form, with
an exponent depending on the fixed number of rims. These are averaged
bounds under the prior kernel. They are not bounds conditional on each
completed stopping history containing colors. The same observation after
exposing a wired layer uses its fresh product of flat kernels, and the
bulk wired assertion handles the initial layer itself.

We next turn the circuit bound into an estimate for a whole collar.
A \emph{window} consists of a working rectangle $R$, a fixed intermediate
enlargement $R^+$, and a larger rectangular buffer, with positive gaps
between them. A traversal test is a path in $R$ from left to right,
staying a fixed distance from its top and bottom, or a rotated copy of
this test. A local wall piece is a component of the cell region
restricted to $R$; pieces joined only outside $R$ remain distinct
locally. We say the region is \emph{anchored} if every component in
the buffer meeting $R^+$ reaches the buffer boundary inside the cell
region. This holds when the region is connected to an exterior seed
outside the buffer. Pieces that touch are treated as one; otherwise the
medial
lattice convention separates their representative paths by lattice
order. Write $F^\eta=\{z:\dist(z,F)\le\eta\}$.

\begin{lemma}[Uniform collar exclusion]\label{lem:collar}
Fix a finite family of such windows and traversal tests. Let
$F_\delta$ be a prescribed matching cell region of color
$\sigma$, anchored in each test. Use one of the following conditional
laws: one actual flat chamber; one selected accessible component facing
at most $k$ connected prescribed sources, with all completed frontiers
of color $\sigma$; or the product class with a common fill. In the second
case
$k$ is fixed and any partition change must satisfy the premise of
\eqref{eq:spatial-wire-density}. Test buffers contain no other pins.
There is no bound on the number of local pieces of one source.

Consider paths in the full underlying dual graph, whether or not their
edges are $N$-open. Let $\mathcal X_\delta(\eta)$ be the event that such
a path makes one of the prescribed traversals, lies in
$F_\delta^\eta$, and is disjoint from $F_\delta$ and from every
augmented matching circuit in the buffer which is attached to
$F_\delta$. Then, uniformly over the stated walls and positive
probability allowed histories,
\begin{equation}\label{eq:spatial-collar-limit}
 \lim_{\eta\downarrow0}\limsup_{\delta\downarrow0}
       \sup \P[\mathcal X_\delta(\eta)]=0.
\end{equation}
An exact cut history may contain colors on its removed side; the
residual search law has the meaning specified above. It is never
conditioned on arbitrary interior $d$ colors.

The same estimate applies simultaneously to paths constrained to lie
in any closed disk bounded by an actual opposite-color $W$ loop, if
that disk is disjoint from $F_\delta$. Such a path automatically avoids
every matching circuit attached to $F_\delta$. The windows and their
buffers may all be chosen inside an open annular band. In particular,
\eqref{eq:spatial-collar-limit} holds along every deterministic
sequence $\eta_\delta\downarrow0$.
\end{lemma}

\begin{proof}
These tests detect every connected passage of a fixed positive
diameter $b$ in a fixed box. Choose an inner square of side a small
fixed multiple of $b$ from a fixed grid around one endpoint, and
stop the connecting path on its first exit from a larger square
of diameter less than $b/2$. The exit side fixes an orientation, and
a fixed enlargement of the vertical extent gives a traversal with
the required margin. There are finitely many grid squares and
orientations. Their working rectangles, intermediate enlargements,
and buffers are all chosen before applying the lemma; anchoring is
required on those intermediate enlargements. The same construction
inside an open band detects a full traversal of the band.

It suffices to prove one left-to-right test. Use the crossing segment
from the last left-side hit before its first right-side hit, and then
erase loops. The result is a simple proper crosscut: its interior
misses both vertical sides.
This trimming is needed because a simple path that returns to its
starting side can leave a third pocket. A proper crosscut has just a
lower and an upper side. Every local obstacle piece disjoint from it is
assigned to one of those sides. Include the top and bottom as virtual
barriers, which the traversal margin already forbids.
Here the pieces are components after restriction to this working
rectangle. Connections outside it retain their wire labels in the
law but do not identify pieces on opposite sides of an endpoint.

We first reduce these assignments, except for a uniformly small event,
to a fixed finite number of possibilities. An anchored piece meeting
a narrower central rectangle contains a path to the boundary of the
working rectangle: take the part of an anchored path up to its first
exit. A piece reaching the top or bottom is already assigned to that
barrier. For a piece reaching the left, in the fixed-width left end strip,
retain a segment from its last hit of the outer vertical to its first
hit of the inner vertical and erase loops. It is a proper crosscut of
that strip. Make a deterministic choice if a piece has
several such paths. Representatives of distinct pieces are disjoint
and can be ordered from bottom to top. The analogous construction is
used at the right side and in the rotated windows. Connections of a
piece between different
sides later impose consistency; they add no choices.

Here is the quantitative reduction in one end strip. Choose a middle
horizontal interval $I$ of length $\ell>0$ and take $\rho>0$ below
both its horizontal clearance from the vertical ends and the fixed
clearance of $R$ inside the larger window buffer. For consecutive
representative arms, let $g_j(x)$ be the difference of their topmost
intersection heights with the vertical at $x$, and put
$d_j=\ell^{-1}\int_I g_j(x)\,\dd x$. These heights have the same
order as the arms: the upper crosscut separates the lower one from the
top of the strip, so a vertical ray from the top to the topmost point
of the lower one first meets the upper one. The arms need not be
graphs. In particular,
\begin{equation}\label{eq:spatial-gap-sum}
                  \sum_j d_j\le H,
\end{equation}
where $H$ is the fixed strip height. The set
$E_j=\{x\in I:g_j(x)\le2d_j\}$ has length at least $\ell/2$.

At lattice scale the separation convention gives $d_j\ge a_0\delta$
for a fixed $a_0>0$. Put $h_j=d_j+\delta\le\beta d_j$, and choose a
fixed large $K$ so that $2Kh_j$ exceeds the cutoff in
\cref{prop:circuits}. If $4K\beta d_j<\rho/2$, a maximal selection of
points in $E_j$ at horizontal separation greater than $16Kh_j$ has
at least $\ell/(64K\beta d_j)$ points, up to a fixed endpoint loss.
Center a disk at the midpoint of the two arm intersections at each
selected point. Both arms meet its radius-$2Kh_j$ disk and leave its
radius-$4Kh_j$ disk, since they continue to the strip ends. The
radius-$8Kh_j$ buffers are disjoint by construction. Since
$8Kh_j<\rho$, they stay inside the allowed window buffer and lie
between the end strip's vertical sides. Every augmented matching
circuit in the intervening annulus meets both proper arms. Orient such
a circuit, stop at its first upper-arm hit after a chosen lower-arm hit,
and retain the arc from
its last preceding lower-arm hit. Its interior misses both arms.
The two proper arms divide the infinite vertical strip into lower,
middle, and upper components. Only the closure of the middle component
meets both arms, so the arc's connected interior lies in that component.
Its closure is contained in the finite end strip. The retained arc
therefore joins the arms geometrically inside the end strip, even if
the rest of the circuit extends beyond its horizontal sides. In
particular, a remote wire alone is not counted as the joining connection.

Successive conditional circuit tests in the completed graph, followed
when needed by the single density comparison
\eqref{eq:spatial-wire-density}, now give
\begin{equation}\label{eq:spatial-gap-failure}
 \P[\text{the two arms have no such local joining circuit}]
                  \le C_k\exp(-c/d_j).
\end{equation}
Here and below $C_k$ is the fixed density factor for the allowed
completion; it is an absolute constant for either exact completion,
including the product class with any number of chambers.
Gaps too large for the preceding packing are simply retained. The
locations of all tests depend only on the prescribed obstacles, not
on the sought traversal. If a matching circuit joins two arms, an
avoiding crosscut cannot assign them to opposite sides.

For any $u>0$, the total probability lost by discarding the gaps
$d_j\le u$ is at most
\begin{equation}\label{eq:spatial-gap-union}
 \sum_{d_j\le u} C_k e^{-c/d_j}
 \le C_k\left(\sum_jd_j\right)
          \sup_{0<d\le u}\frac{e^{-c/d}}{d}
 \longrightarrow0\qquad(u\downarrow0).
\end{equation}
At most $H/u$ larger gaps remain, together with a bounded number
adjacent to virtual barriers. A proper traversal can split the
ordered representatives at only one of these ranks. Repeating this
on the finitely many sides leaves a finite product of deterministic
rank lists, bounded in terms of $u$ and the windows. A piece joined
around another side must receive the same assignment there, which
only removes rank choices. Trials for different gaps may overlap;
the union bound in \eqref{eq:spatial-gap-union} does not require their
independence.

\begin{figure}[t]
\centering
\begin{tikzpicture}[x=0.8cm,y=0.8cm,>=Latex]
 \draw[gray] (0,0) rectangle (6,4);
 \draw[densely dashed,gray] (1,0)--(1,4) (5,0)--(5,4);
 \draw[very thick,teal!70!black] plot[smooth] coordinates
   {(0,0.8) (1,1.15) (2,0.9) (3,1.35) (4,1.0) (5,1.2) (6,0.9)};
 \draw[very thick,teal!70!black] plot[smooth] coordinates
   {(0,2.7) (1,2.35) (2,2.55) (3,2.1) (4,2.45) (5,2.3) (6,2.7)};
 \draw[teal!60!black,densely dashed] (2,1.72) circle (0.65);
 \draw[teal!60!black,densely dashed] (4.1,1.73) circle (0.65);
 \draw[<->] (3,1.38)--node[right,fill=white,inner sep=1pt] {$g_j(x)$}
                  (3,2.08);
 \node at (3,-0.35) {local annuli join ordered arms};
 \begin{scope}[xshift=7.5cm]
  \draw[gray] (0,0) rectangle (6,4);
  \draw[very thick,teal!70!black] plot[smooth] coordinates
    {(0,0.6) (1,0.9) (2,0.75) (3,1.05) (4,0.8) (6,0.65)};
  \draw[very thick,teal!70!black] plot[smooth] coordinates
    {(0,3.45) (1,3.1) (2,3.3) (3,3.0) (4,3.25) (6,3.4)};
  \draw[very thick,orange!80!black] (3,1.42)--(3,2.63);
  \draw[teal!60!black,densely dashed] (3,1.05) circle (0.5);
  \draw[teal!60!black,densely dashed] (3,3.0) circle (0.5);
  \draw[blue!65!black,thick,->] (0,2.0).. controls (1.0,1.7) and
    (1.7,2.25)..(2.45,2.05);
  \draw[blue!65!black,thick] (3.55,2.05).. controls (4.3,1.85) and
    (5,2.2)..(6,2.0);
  \node[anchor=south west,inner sep=1pt] at (3.5,2.18) {$K_m$};
  \node at (3,-0.35) {a lower--upper connector blocks a crosscut};
 \end{scope}
\end{tikzpicture}
\caption{The two geometric steps in collar exclusion, drawn
schematically. Representative arms can double back; their topmost
vertical encounters remain ordered. The dashed matching circuits are
local geometric circuits, with any artificial part contained in the
forbidden cell region. In the right panel the compact gap segment
$K_m$ is forbidden to a sufficiently thin collar.}
\label{fig:spatial-collar}
\end{figure}

We finish by a compactness argument, keeping its choices deterministic.
If \eqref{eq:spatial-collar-limit} failed, there would be deterministic
allowed histories and walls at meshes $\delta_n\downarrow0$, widths
$\eta_n\downarrow0$, and a fixed positive lower bound for the traversal
probability. First fix $u$ so that \eqref{eq:spatial-gap-union}, summed
over the sides, is less than one quarter of that bound. One of the
finitely many surviving rank choices then has a fixed positive event
mass along a subsequence. We do not condition the law on that event.
Take Hausdorff limits, in the central buffered rectangle, of the
unions assigned below and above, including the virtual barriers.

Suppose first that the two limits touch there. Unless a virtual
barrier is involved, there are deterministic points on two
oppositely assigned pieces at distance $r_n\downarrow0$. Anchoring
makes both pieces cross the annuli from scale $r_n+O(\delta_n)$ to
a fixed outer scale. Any augmented matching circuit in such an
annulus joins them. By \eqref{eq:spatial-attached-arm}, the probability
of a traversal with this assignment tends to zero. If one limiting
piece is a virtual barrier, the fixed traversal margin rules out
the opposite assignment: a connected piece reaching that margin
lies on the barrier's side of the proper crosscut.

It remains that the two closed limits are disjoint in the central
rectangle. On a central vertical line choose a consecutive lower
encounter $a$ and upper encounter $b$. Such a transition exists
because the virtual bottom and top have those assignments. The open
segment $(a,b)$ misses both limiting unions. Fix an integer $m$.
At each actual endpoint choose $m$ separated circuit radii, all less
than $|a-b|/8$ and the available buffer clearance; call their smallest
radius $r_m>0$. If both endpoints are virtual, take any fixed
sufficiently small $r_m$ in their traversal margins instead.
Only after this choice, trim the gap by less than $r_m/2$ at its
ends and call the remaining compact segment $K_m$.
It has positive distance from the whole limiting obstacle union, not
just from its intersections with the vertical. This clearance holds
for the discrete obstacles for all large $n$, and then
$K_m\cap F_{\delta_n}^{\eta_n}=\varnothing$.

Center the endpoint tests at discrete points converging to $a$ and
$b$. An attached circuit at any of the chosen radii meets $K_m$:
it surrounds its center, and the ray into the gap crosses the circuit
between that radius and twice that radius. Its other intersection is
with the endpoint's anchored piece. For a virtual endpoint the short
segment in the fixed traversal margin supplies the connection instead.
Consequently a successful endpoint circuit at each actual endpoint,
together with $K_m$, is a connected set from a lower-assigned piece
or barrier to an upper-assigned one, wholly in the rectangle. Every
proper left-to-right crosscut with that assignment meets this set.
It cannot do so: it avoids the pieces and the attached circuits, and
the collar condition excludes $K_m$.

For fixed $m$, the probability that either required endpoint
connection fails has limit superior at most
$2C_k(1-p)^m$, by the separated-scale estimate and, if needed, one
density comparison. Letting $m$ increase contradicts the retained
positive event mass. The order of choices here matters: the finite
endpoint radii precede the trimmed segment, so every requested
circuit really reaches that segment. This proves uniformity in
\eqref{eq:spatial-collar-limit}.

Finally let a closed disk be bounded by an actual $W^{-\sigma}$ loop
and disjoint from $F_\delta$. A connected augmented matching circuit
attached to $F_\delta$ cannot enter this disk. Its artificial edges
are in $F_\delta$, outside the disk, and its actual $W^\sigma$ edges
cannot cross or share a vertex with the opposite-color boundary.
A dual path inside the disk therefore automatically avoids all those
circuits. The existence event already ranges over every underlying
dual path, so the bound holds simultaneously over the possible disks,
without conditioning on a future choice of one. Choosing all the
windows with fixed margin inside an open band gives the last
localization assertion. Monotonicity in $\eta$ gives the conclusion
for $\eta_\delta\downarrow0$.
\end{proof}

\subsection{Counts and comparison of complete patterns}

\begin{proposition}[Cluster and transmission counts]
\label{prop:cluster-counts}
The following bounds hold in fixed buffered boxes, uniformly in the
mesh and in allowed conditional histories. A joint count of actual
clusters across chambers in items~1 and~3 requires the product class
with a common fill. Item~2 has the interior scope stated there.
\begin{enumerate}
\item For fixed $b>0$, the number of unsearched free $N$ clusters
having a connected passage of diameter at least $b$ in a fixed box
has an exponential tail. Its constants depend on $b$ and the fixed
box geometry, not on the roughness of the surrounding walls or on
the number of chambers in the realization with a common fill.
\item For one chamber, let a rectangle and its joining buffer lie
wholly in the unprescribed accessible interior. Use the relative dual
restriction with separate side ports defined in the proof. The number of distinct
components of the restricted $N$ graph crossing between two specified
opposite sides has a geometric tail. Consequently, if
$Q_0\Subset Q_1$, the number of restricted components in $Q_1$
which meet $Q_0$ and a vertex incident to an edge outside the
restriction has a geometric tail, provided the fixed covering
rectangles in $Q_1\setminus Q_0$ and every joining buffer have the
same unprescribed interior clearance. Arbitrary admissible $(s,W)$
prescriptions outside a larger buffer are allowed. The same interior
assertion holds in the original wired law before removal of its
boundary cluster.
\item Counts of distinct free $N$ clusters of the residual kernels
enclosing or passing through a fixed
aspect-ratio annular shell have moments of every fixed order, uniformly
after rescaling the shell down to lattice order. Under a stopping
history these are the unsearched clusters. A distinguished wired
cluster adds at most one. For arbitrary
independent chambers with pairwise disjoint interiors this last
statement is available separately in each chamber; it asserts no
bound for their total count without a common fill. The precise
no-shadowing consequence at indexed discoveries is stated in
\cref{cor:indexed-no-shadowing} below.
\end{enumerate}
\end{proposition}

\begin{proof}
For item~1, start with the connected cell region of the given
exact cut history; for a single flat chamber use its exterior side.
Explore with a deterministic connected seed attached to that region.
Add one seed vertex at a time and immediately reveal the entire
previously undiscovered $N$ cluster it meets. The seed eventually
visits every vertex of the box. Record only discoveries with the
specified passage, choosing one such passage by a deterministic rule.
At every recorded stop, the initial region, seed, and all clusters
already discovered form one connected cell region. By
\cref{lem:finite-cuts}, the residual laws are exactly the flat
kernels, jointly in the class with a common fill when there are several
chambers.

A finite family of smaller collar windows detects every passage
of diameter at least $b$. Freeze the cell region $F_i$ at the
$i$th recorded stop, with $F_0$ the initial region, before continuing
the search. An actual unsearched $N$ path is disjoint from $F_i$ and
cannot cross the actual part of an augmented matching circuit.
Thus \cref{lem:collar} gives an $\eta>0$ such that, conditionally on
any such frozen history, the probability that \emph{some} still
unsearched qualifying passage lies wholly in $F_i^\eta$ is at most
a fixed $p<1$. This event is existential over future passages; no
future location or visibility is conditioned upon. It also bounds
the event that the next recorded cluster exists and its selected
passage lies wholly in $F_i^\eta$.

On the complementary event, the next passage contains a point at
distance greater than $\eta$ from $F_i$. Such points from successive
records are mutually $\eta$-separated, because each earlier cluster
is included in every later frozen region. A bounded box contains
at most $M=M(\eta)$ such points. There is no cost for the number
of intervening microscopic discoveries: they only enlarge the same
connected region. In the product case reflect the residual chamber
spins to one matching color and use the exact realization with a common fill
for this single existential collar event. In particular, no union
bound over the number of holes is taken. A test buffer meeting the
initial wall is covered by the same augmented graph.

If at least $n$ records occur, at most $M$ of the first $n$ can be
of the separated-point kind. Conditional iteration for the other
indices, including the initial stop, gives
\begin{equation}\label{eq:spatial-count-tail}
 \P[\text{at least }n\text{ recorded discoveries}]
       \le \sum_{j=0}^{M}\binom{n}{j}p^{\,n-j}
       \le C e^{-c n}.
\end{equation}
The same proof starting from a later allowed stop bounds its
unsearched count. It does not bound the number already fixed
by that history. Meshes above the small threshold used by the collar
lemma are absorbed by the deterministic number of vertices in the
fixed box.

For item~2 we adapt the crossing exploration underlying
\cite[Section~19.1, especially Lemmas~19.1 and~19.4]{DKLM2026}.
First fix the relative planar dual of the given restricted $N$
lattice quadrilateral. Its medial boundary follows the four sides
of the rectangle. In the auxiliary $W$ graph, truncate a dual edge
at that boundary when it exits, and give each boundary incidence
its own degree-one port. There are no edges in an exterior side
column and no identifications between distinct ports. Every retained
edge or half-edge inherits the state of its actual $W$ edge. This is the matched dual
restriction used for every $W$ connectivity test below; its
placement differs from the geometric rectangle by $O(\delta)$,
which is absorbed by the joining buffer.

If $J_N$ distinct restricted $N$ components cross the quadrilateral,
choose one simple crossing in each, ordered from bottom to top.
Between each two successive crossings, relative planar duality gives
a parallel open crossing of this matched $W$ graph. The separators
in two successive gaps belong to different components of that graph:
a connection inside it would cross the intervening $N$ path, and
the separate ports give no route around an endpoint. Consequently
\[
                  J_N\le J_W+1,
\]
where $J_W$ counts crossing components in precisely the matched
$W$ graph. The probabilistic search uses only these $W$ crossings;
it never certifies a separator against an unrevealed upper $N$ crossing.

Distinct $W$ crossing components have a consistent vertical order:
one component cannot have crossing paths on both sides of another,
because its joining path would cross the latter. In each component
choose a lowest simple proper crossing, with a deterministic tie
rule; the lower boundary of the union of its crossing paths supplies
such an open crossing. Enumerate these representatives from below,
with all connectivity restricted to the matched graph.
For a candidate $\gamma$, the event that it is the next distinct
crossing is determined by the $(s,W)$ configuration in its closed
below envelope. All lower candidate crossings lie there, and any
connection from a lower component to $\gamma$ first meets $\gamma$
from below. A proposed connection between lower pieces through the
region above $\gamma$ would first join $\gamma$ as well. This is
the stopping-line property for the enumeration.

After selecting $\gamma$, reveal the whole below envelope, with
endpoint spins before edge states. Calling the $W$ parity even,
extend endpoint half-edges only for the following geometric closure,
and close the envelope to an even domain $D_\gamma$ through the
unused side and bottom buffer outside the
central test balls. All earlier queried regions and their possibly
unfavorable frontiers are now inside $D_\gamma$; its boundary
segments in the central balls are exactly matching segments of
$\gamma$. The endpoint extensions are not edges of the counted
matched graph. The transcript is a cylinder measurable in $D_\gamma$.
Equivalently, one can disintegrate over that full cylinder after
the selection. The envelopes are nested as the search moves upward,
so the next transcript still has this form.

In a central vertical strip with a buffer continuing above the
rectangle, use the partial-circuit construction of
\cite[Lemma~19.4 and the proof of Lemma~19.1]{DKLM2026}. That lemma
is stated for the plane: it allows arbitrary same-parity data in
revealed even domains whose in-ball frontier is matching, and tests
a circuit in the matching graph augmented by those domains.
To transfer its bound to a finite conditional kernel, choose a
complete deterministic outer circuit $\Gamma_{\rm out}$ of the $W$
lattice, with its incident edge layer inside the unused buffer and
outside the test ball. Let $G^\circ$ be the graph induced by the $W$
vertices strictly inside $\Gamma_{\rm out}$, and let $V_\partial$ be
its vertices adjacent to that circuit. Analytically disintegrate the
old conditional law over the spins at $V_\partial$ and a complete
exterior and interface realization, fixing all endpoint spins before
edge states. For each fixed contact-spin vector, the positive exterior
fibers give laws on $G^\circ$ with identical retained interior pins
and local edge factors; only the same-color
partition of $V_\partial$ induced from outside varies. The interface
factors are constants once these endpoint spins are fixed. Every
exterior connection reaches $G^\circ$ through $V_\partial$. This
includes inherited abstract wires, whose terminals are outside the
unprescribed buffer.

For a sought plus circuit, replace this exterior partition by one
class on all minus contacts and singleton exterior classes on plus
contacts, retaining every interior pin and wire. This is the adverse
ordered boundary condition of \cref{prop:comparison}: it preserves or
merges every old minus connection and removes only plus exterior
connections. Realize it by pinning the vertices of $\Gamma_{\rm out}$
minus and its edges open. Pin each inward incidence $vc$, with
$v\in V_\partial$ and $c\in\Gamma_{\rm out}$, open minus if $s_v=-$
and closed if $s_v=+$; pin any interior chord of $\Gamma_{\rm out}$
open minus as well. These are compatible states after the endpoint
spins are fixed. An old minus-open incidence remains open, an old
plus-open incidence is closed, and an old closed incidence is either
kept closed or opened minus. The comparison is between kernels on
the common graph $G^\circ$; it does not impose the old exterior pins
on this new realization. Reflect colors for a sought minus circuit.

Under the complete cut $\Gamma_{\rm out}$, the finite interior
$(s,W)$ law and the plane law conditioned on the same cut have the
same flat kernel. Condition both on the identical retained interior
pins, contact spins, inward incidence states, and interior chords.
Cover the full added interface layer by finitely many small even
domains whose closures remain outside the test ball. Together with
$D_\gamma$, they make this conditioning measurable in the auxiliary
domain union, while adding no in-ball frontier segment or augmented
edge in the tested annulus. Lemma~19.4 applies in this plane
representation, with exactly the matching in-ball frontier of
$\gamma$, and gives the uniform positive partial-circuit bound.
The ordered comparison holds for every old fiber, so averaging with
its original probability introduces no factor. The original wired law
uses the same $(s,W)$ comparison before the cut. Multiple frontier
segments are handled by the finite-family conclusion of Lemma~19.4;
the polar Markov statement in Lemma~14.13 itself has only two arcs.

Intersect a fixed chain of these tests under the conditional pin
FKG law. With a fixed positive probability their augmented tube
joins the filled lower part to above the rectangle. It stays in
the central strip, where its last exit from $D_\gamma$ is on
$\gamma$ itself. The remaining connection is actual and lies in
the interior of the matched graph up to the top. A later parallel
crossing of the same color would meet it and join $\gamma$ within
that graph, while a crossing of the opposite color cannot cross it.
Thus the tube event rules out another distinct $W$ crossing.

These tube tests are used to bound the conditional probability;
their upper data are not added to the next search transcript.
For every transcript at the $i$th representative,
the conditional chance of a further one is at most $1-p$.
Iteration gives a geometric tail for $J_W$ and hence for $J_N$.
Connections outside the rectangle play no role because both counts
use restricted components. The original wired law has the same
interior $(s,W)$ comparison by \cref{prop:circuits}.

A fixed finite rectangle cover of the shell between $Q_0$ and $Q_1$
assigns each transmitting restricted component to a crossing of one
of these rectangles. The small lattice layer in the incidence
definition is absorbed by the fixed clearance. This proves the
transmission assertion.

For the shell assertion, the clusters counted are unsearched global
free clusters of the residual law, rather than distinct restricted
components. At shell
radius $r$, a fixed finite cover assigns every enclosing or passing
cluster a connected passage of diameter at least $c_2r$ in one
of its boxes. Apply item~1 after rescaling by $r$. Its collar proof
is uniform when $\delta/r$ is small, including boxes met by matching
walls in the realization with a common fill. When $\delta/r$ is bounded
below there are only a bounded number of vertices in the rescaled
boxes. The resulting exponential tail gives every fixed moment,
uniformly down to lattice order. For disjoint independent chambers
apply this assertion separately in each one. A pair of vertices in
different chambers shares no cluster, which is the only
cross-chamber fact needed for the later moment estimates.
\end{proof}

\begin{corollary}[Counts in an enclosing box]
\label{cor:global-cluster-counts}
Fix a bounded box $L$ and $b>0$. Suppose the finite lattice domain
and the supports under consideration lie in $L$, and place $L$ in
a fixed larger regular ambient box. For one free flat $N$ law, or
the residual product law with a common fill conditional on its initial
exact cut data, let
\[
 \mathcal N_\delta(b)
   =\#\{\,\text{free }N\text{ clusters }C:\diam C\ge b\,\}.
\]
There are constants $C,c>0$, depending on $L,b,t$ but not on the
initial walls or the mesh, such that
\begin{equation}\label{eq:spatial-global-count-tail}
       \P[\mathcal N_\delta(b)\ge n\mid\text{initial cut data}]
                       \le C e^{-cn}.
\end{equation}
The unconditioned free dual current is included. For the original
wired current, expose its boundary cluster first; its free
descendants satisfy this estimate conditionally, and adding the
distinguished boundary cluster increases the count by at most one.
The boxes and all cover buffers are fixed before the mesh limit.

There is also a pathwise consequence for a seed exploration. Suppose
a compact set $Q$ is at distance at least $r>0$ from the seed and
initial removed support. List the discovered free clusters whose
support meets $Q$, or whose union of odd holes meets $Q$. If
$\mathcal L_\delta(\tau;Q,r)$ is the number listed up to any stop,
then, for all sufficiently small $\delta$,
\begin{equation}\label{eq:spatial-contributor-domination}
       \mathcal L_\delta(\tau;Q,r)
                       \le \mathcal N_{\delta,\mathrm{initial}}(r/3).
\end{equation}
Add one on the right if the distinguished wired cluster is listed.
This is domination by the initial count, not a tail conditional on
the completed stop.
\end{corollary}

\begin{proof}
Cover $L$ by finitely many inner squares of side a small fixed
multiple of $b$. Around each choose a passage square and a larger
buffer, with fixed ratios, in the regular ambient box. If a
connected cluster has diameter at least $b$, choose two of its
points at that separation and an inner square containing the first.
The path to the second has a first exit from the passage square;
its initial portion is a connected passage of diameter at least
$c_1b$ there, for a fixed $c_1>0$. Let
$\mathcal N_{\delta,i}^{\rm passage}(c_1b)$ be the item~1 count
of unsearched global free clusters of the initial residual law
having such a passage in the $i$th square.
Every global cluster is counted in at least one of these counts.
Consequently
\[
        \mathcal N_\delta(b)
            \le \sum_{\text{cover squares }i}
                 \mathcal N_{\delta,i}^{\rm passage}(c_1b).
\]
The cover is fixed before $\delta$. Its passage buffers may cross
the original boundary or the revealed wired layer: extend the
initial cell region to the ambient box and use the augmented
matching completion of \cref{prop:circuits}. Item~1 of
\cref{prop:cluster-counts} and a finite union bound give
\eqref{eq:spatial-global-count-tail}. For the original free dual
law the matching primal boundary is the initial fill. Behind the
wired $N$ boundary cluster, \cref{lem:finite-cuts} gives precisely
the product law whose common initial fill is that connected cluster
joined to the exterior.

For \eqref{eq:spatial-contributor-domination}, every discovered
cluster meets the seed. If its support meets $Q$, it has diameter
at least $r-O(\delta)$. If instead one of its odd holes meets $Q$,
that hole lies in the convex hull of the cluster's cell region:
a bounded complementary component is contained in that convex
hull. A point of the hole is therefore at distance at most
$\diam C+O(\delta)$ from the seed contact on $C$. This again gives
$\diam C\ge r-O(\delta)$. The listed clusters are distinct members
of the initial free collection, proving the pathwise inequality.
This argument covers a cluster which surrounds $Q$ without meeting it.
\end{proof}

\begin{lemma}[Distance to a chamber-side rim]\label{lem:chamber-rim-distance}
Let $U_\delta$ be a resolved residual chamber of an exact cut
exploration, let $S_\delta$ be the revealed support including the
initial boundary and seed as needed, and let $\Gamma_\delta$ be its
own chamber-side
$W$ perimeter. For any subset $C_\delta$ of that chamber,
\begin{equation}\label{eq:spatial-own-rim-distance}
 \sup_{x\in C_\delta}\dist(x,\Gamma_\delta)
       \le \sup_{x\in C_\delta}\dist(x,S_\delta)+O(\delta).
\end{equation}
The constant is the lattice resolution constant. In particular, if
$C_\delta\to C_*$ and $S_\delta\to S_*$ in Hausdorff distance with
$C_*\subset S_*$, the left side tends to zero.
\end{lemma}

\begin{proof}
For $x\in U_\delta$, take a segment to a nearest point of
$S_\delta$. Its first exit from the resolved chamber meets that
chamber's own medial perimeter, within $O(\delta)$ of the
corresponding $W$ edge. The distance to this first exit is at most
the length of the segment. Taking suprema proves the inequality.
This uses the own perimeter even if the nearest support point
belongs to another nearby part of a pinched wired layer. For the
original free chamber the same argument uses its polygonal boundary
as $S_\delta$. The Hausdorff conclusion follows from the one-sided
inclusion $C_*\subset S_*$.
\end{proof}

\begin{corollary}[No shadowing at recorded stops]
\label{cor:indexed-no-shadowing}
Fix a buffered box $Q$ and $b>0$. At a stop $\tau$ of an
exterior-attached seed exploration, let $F_\tau$ be its frozen
connected cell region, including the initial walls, all visited seed
vertices, and all clusters already revealed. Let
\[
 \Xi_\delta(\tau;b,\eta)
  =\left\{\begin{array}{c}
   \text{some unsearched }N\text{ cluster has a connected}\\[-2pt]
   \text{passage }\gamma\subset Q,\quad \diam\gamma\ge b,\quad
              \gamma\subset F_\tau^\eta
          \end{array}\right\}.
\]
For allowed exact cut or $(s,W)$ stopping histories $\mathcal H_\tau$
with the supported residual product class,
\begin{equation}\label{eq:spatial-frozen-shadow}
 \P[\Xi_\delta(\tau;b,\eta)\mid\mathcal H_\tau]
       \le\omega_{Q,b}(\eta,\delta),\qquad
 \lim_{\eta\downarrow0}\limsup_{\delta\downarrow0}
       \omega_{Q,b}(\eta,\delta)=0,
\end{equation}
uniformly over positive probability histories. Exterior colors
rendered irrelevant by a completed cut are allowed. Arbitrary
interior $d$ pins, or conditioning on a later visibility event,
are not part of this statement.

For an exploration in the enclosing box of
\cref{cor:global-cluster-counts}, index stops by their rank among
discoveries of clusters of diameter at least $b$, and let
$\tau_0$ be the initial stop. All intervening smaller discoveries
are included in each frozen $F_{\tau_i}$. If $D_\delta(b)$ is the
number of these records, then, for every integer $K$,
\begin{equation}\label{eq:spatial-indexed-shadow}
 \P[\exists\,0\le i\le D_\delta(b):
             \Xi_\delta(\tau_i;b,\eta)
       \mid\text{initial cut data}]
 \le C e^{-cK}+(K+1)\omega_{Q,b}(\eta,\delta).
\end{equation}
The indices here are recorded macroscopic ranks, not raw discovery
numbers.

If a random width $\eta_\delta$ tends to zero only on a limiting
event $E$, the usable deterministic inclusion is, for fixed $\eta$,
\begin{equation}\label{eq:spatial-random-width}
 E\cap\Xi_\delta(\tau;b,\eta_\delta)
 \subseteq
 \Xi_\delta(\tau;b,\eta)\ \cup\
       \bigl(E\cap\{\eta_\delta>\eta\}\bigr).
\end{equation}
Thus one first fixes $\eta$, takes the mesh limit (and Fatou's
inequality when passing to a limit event), and then sends $\eta$ to
zero. No estimate is asserted after conditioning on $E$.
\end{corollary}

\begin{proof}
An actual unsearched $N$ passage avoids $F_\tau$ and the actual
edges of every attached augmented matching circuit. A finite
collar-window cover of $Q$ therefore gives
\eqref{eq:spatial-frozen-shadow} from \cref{lem:collar}, including
windows met by the initial wall. The exact cut factorization
identifies the residual marginal for every permitted history.
The bound is for the existential future event under that marginal;
it does not reveal a particular future passage.

On $D_\delta(b)\le K$, there are at most $K+1$ frozen indexed
stops. Apply \eqref{eq:spatial-frozen-shadow} at each one and use
a union bound. The complementary event is bounded by
\eqref{eq:spatial-global-count-tail}, which yields
\eqref{eq:spatial-indexed-shadow}. Finally,
\eqref{eq:spatial-random-width} follows from monotonicity of
neighborhoods. It gives the stated order of limits even when the
width was chosen from a future Hausdorff containment.
\end{proof}

For subsequent analytic arguments an additive mixing estimate is
not sufficient: the prescribed microscopic patterns may themselves
have very small probability. We therefore record a multiplicative
comparison. A \emph{complete pattern} in a finite vertex-and-edge
set specifies all its $s$ spins and all its $W$ edge states. Every
specified edge includes its endpoints.

\begin{proposition}[Pattern comparison]
\label{prop:pattern-comparison}
Fix regular compact sets $Q_i$ and filled buffered neighborhoods
$B_i$, $1\le i\le m$, with the $B_i$ pairwise disjoint and a fixed
positive separation between $Q_i$ and $\partial B_i$. Assume each
$B_i$ is contained in the domain, with no prescribed data there
other than the pattern being tested in $Q_i$. Exterior data may be
arbitrary admissible $(s,W)$ data and their retained connections.
There is a constant $C$ such that, for complete patterns $p_i$,
\begin{equation}\label{eq:spatial-product-patterns}
 C^{-m}\prod_{i=1}^m\pi_i(p_i)
 \le \mu[p_i\text{ on every }Q_i]
 \le C^m\prod_{i=1}^m\pi_i(p_i).
\end{equation}
Here $\mu$ is any such conditional flat law, and $\pi_i$ is the
plane pattern law in $Q_i$. The constant depends on the fixed
buffer geometry and $t$, not on the patterns or on their number
of microscopic prescriptions. Copies under parity-preserving lattice
translations and reflections satisfy the same assertion.

In particular, the single-box pattern laws have uniformly bounded
Radon--Nikodym density ratios, and the same remains true locally
after conditioning on any collection of separated patterns in the
other boxes.
\end{proposition}

\begin{proof}
Consider first a single $Q$ inside its buffer $B$. Prescribe a
complete admissible pattern in four batches: its plus vertices,
its minus vertices, its increasing compatible edge prescriptions,
and its decreasing compatible edge prescriptions. After the first
two batches every edge endpoint is fixed. A disagreeing edge is
already determined closed. On monochromatic edges the two remaining
batches mean open plus or closed minus, and closed plus or open
minus, respectively. Each batch is monotone in one of the two
orders of \cref{prop:comparison}. The number of batches is four,
independently of the pattern size.

Fix the preceding batches and consider an increasing current batch
$E$. Choose two deterministic annular buffers between $Q$ and
$\partial B$, and fix a plus circuit on the farther one. Let $\lambda$
be the law inside that deterministic circuit with the preceding
prescriptions retained. Comparison gives
\begin{equation}\label{eq:spatial-batch-upper}
                         \mu(E)\le\lambda(E).
\end{equation}
In the nearer annulus, conditional on the preceding prescriptions,
a random plus circuit has probability at least $p>0$. To see this
even if those prescriptions contain both signs, cover that annulus
by buffered test boxes missing $Q$; the exterior bound for one parity
and FKG supply its circuit with the same positive constant.
Search outside-in, inspecting nothing on the sealed side, and stop
at the first such circuit in a fixed exploration. Its interior
law is plus wired by \cref{lem:finite-cuts}, and hence dominates
the restriction of $\lambda$, with the same preceding prescriptions.
Consequently
\begin{equation}\label{eq:spatial-batch-lower}
                         \mu(E)\ge p\lambda(E).
\end{equation}
The decreasing batch uses minus circuits and the reversed order.

The law $\lambda$ is common to any two exterior conditions under
comparison. Applying \eqref{eq:spatial-batch-upper}--
\eqref{eq:spatial-batch-lower} to each of the four conditional batch
probabilities yields bounded ratios for the complete pattern
probabilities. Zero-probability patterns are exactly the common
local incompatibilities and are harmless. These constants remain
valid when the comparison law is a flat exhaustion of the plane.
The common plane spin/percolation law obtained from flat exhaustions in
\cite[Theorem~17.1]{DKLM2026} therefore gives the comparison with
$\pi_i$ by taking that exhaustion limit. For fixed $W$ parity this
law is invariant under parity-preserving lattice translations and
reflections, with the compatible global residue action. An
automorphism exchanging the lattices acts also by exchanging which
parity is called $W$; it is not invariance of the same fixed edge
decoration under that automorphism. The height-gradient marginal
has the full corner-grid symmetries used later.

For several boxes, condition successively on the patterns in the
other boxes. These constraints lie outside the current filled
buffer, so the single-box estimate applies with unchanged constants.
The chain rule gives \eqref{eq:spatial-product-patterns}. Conditioning
on separated patterns and repeating the same argument proves the
last assertion. Notice that the proof is multiplicative at every
stage; it is not an inference about rare events from a total
variation error.
\end{proof}

\subsection{Extending separation witnesses}

The preceding proposition concerns one spin parity and its edge
decoration. The other spins color complementary components, so their
admissible assignments can depend on connections outside a tested
pattern. We now localize that dependence. In this subsection choose
$W$ on the parity carrying the flat wall. The complementary $N$
components in the chamber are all free. This restriction is needed
only for the color statement here and below; the wired bulk circuit
case of \cref{prop:circuits} remains available.

Fix regular boxes $Q\Subset Q'\Subset B$ with fixed positive gaps.
For a geometric closed set, its complete pattern includes every $W$
edge meeting the set and both endpoint spins; the resulting one-mesh
boundary layer is part of the pattern. Form the restricted $N$ graph
using precisely the dual edges whose $W$ states that pattern
determines. A restricted component is a \emph{terminal} if it meets
$Q$ and contains a vertex incident to an unprescribed $N$ edge.
Every component meeting $Q$ that is not a terminal is sealed by the
pattern, so its color is independent of the exterior. This incidence
definition avoids losing an edge in the lattice boundary layer.
Select one marked vertex of each terminal by a deterministic rule.
A vector of signs on these marks is \emph{compatible} with an
exterior configuration if no two oppositely assigned marks become
connected in the full $N$ graph.

\begin{lemma}[Buffered separation witnesses]
\label{lem:buffered-witnesses}
Fix $p>0$ and an integer $k$. Conditional on any complete pattern
in $Q'$ with at most $k$ terminals, suppose a specified terminal
assignment has compatibility probability at least $p$ in a
$W$-flat domain containing $B$. There is
$c=c(p,k,Q,Q',B,t)>0$ such that, for every sufficiently small mesh,
with conditional probability at least $c$, every opposite pair
of marks is separated by an actual open monochromatic $W$ loop
contained in $B$. Thus compatibility is witnessed without using
the configuration beyond $B$. The estimate is uniform in the
microscopic pattern and in the enclosing domain.

Admissible $(s,W)$ exterior data beyond a fixed larger buffer are
allowed. No transverse $d$ colors are conditioned upon in this
statement.
\end{lemma}

\begin{proof}
In a finite $W$-flat domain, planar cut-cycle duality says that two
marks disconnected in $N$ are separated by a simple monochromatic
$W$ loop, allowing portions of the matching outer perimeter. We
will replace possibly long loops by loops in successively enlarged
cores. Write $\mathcal O$ for the set of required opposite pairs
and $m=|\mathcal O|\le k(k-1)/2$. If $m=0$ there is nothing to prove.
Choose regular closed cores
\[
       Q'=K_0\Subset K_1\Subset\cdots\Subset K_m\Subset B.
\]
In each gap fix a middle annular band and a wider band compactly
inside the gap, with further positive margins. At lattice scale
the core includes every $W$ edge meeting the closed core and its
endpoints. This convention counts a crossing on its boundary; all
such edges of $K_m$ lie in $B$ for small mesh.

For a complete pattern $g$ on $K_j$, let $\mathcal S(g)\subseteq
\mathcal O$ be the pairs separated by some actual monochromatic
loop in that lattice core. This set is determined by $g$ and can
only increase when the core pattern is enlarged in the same
configuration. A pair outside $\mathcal S(g)$ is called missing.
We prove the following one-step statement. If, under the original
law conditional only on $g$, compatibility has probability at
least $a>0$ and some pair is missing, then
\begin{equation}\label{eq:spatial-witness-step}
 \P[\text{compatibility and }|\mathcal S(P_{j+1})|
             >|\mathcal S(g)|\mid P_j=g]\ge r(a,k)>0.
\end{equation}
The constants use the fixed bands of this step and are uniform in
$g$ and the mesh. Here $P_j$ denotes the complete core pattern.

First choose a finite type of witness configuration. On compatibility,
among simple loops separating a missing pair choose one whose enclosed
set $I$ of marks is inclusion-minimal, with a deterministic rule for
ties. It contains a nonempty proper subset of the marks. No missing
pair is wholly inside its disk. To prove this, let $C$ be the selected
loop and suppose a missing pair $a,b$ lies inside. Choose a witness
loop $D$ for that pair. If it meets $C$, their colors agree, since
opposite-color open $W$ paths cannot meet. Consider the finite planar
graph
\[
             H=(C\cup D)\cap\overline{\operatorname{int}C}.
\]
Every underlying dual path from $a$ to $b$ leaving
$\operatorname{int}C$ crosses $C$, and every such path staying
inside crosses $D$. Thus $H$ separates $a$ from $b$ in the full
plane. Planar cut-cycle duality extracts a simple cycle of $H$
separating them. Its disk lies in
$\overline{\operatorname{int}C}$ and contains one of $a,b$ but not
the other, so its marked set is a nonempty proper subset of $I$.
All its edges have the common color. If $C$ and $D$ are disjoint,
the disk of $D$ must lie inside that of $C$ and gives the same
contradiction. This proves the claim, including shared-edge and
tangency cases.

Call $C$ the target. For every missing pair wholly outside $I$,
choose a witness loop and record its color. Record also the target
color and $I$. Reflect the colors for the description so that the
target is plus. A chosen minus witness is disjoint from the target;
its disk is either disjoint from the target disk or contains that
disk. Of its two marked endpoints record the one on the side away
from the target disk. The whole type consists of these data,
including the ordered away-endpoint vector and the fixed mark
$i_0=\min I$. There are at most $L(k)<\infty$ types, independently
of the shapes of all loops. If the compatibility probability is at
least $a$, one type event $\mathcal T$ has probability at least
$a_0=a/L(k)$. Fix that type now. The sampling law remains the
original law conditional on $g$; it is not conditioned on
$\mathcal T$.

From its fixed away marks perform a deterministic exploration of
the full components of the underlying dual graph whose edges do
not cross an open $W^-$ edge. These are components of the
complement of $W^-$, not components of $N$. Explore through already
known core edges as well as new ones, revealing endpoint spins
before each new edge state, until each component is exhausted.
There are at most $k$ distinct source components. The transcript
$h$ is the cylinder of the queried values, so its residual law
$\mu_h$ is an allowed pin law in \cref{prop:comparison}. In
particular it remains positively associated. The interior frontier
of a full source component consists of actual $W^-$ edges; it is
not a frontier artificially stopped at the old core. Original
outer-boundary conditions are retained, and the middle-band test
windows are away from the old core and the outer boundary.

Let $F_h^0$ be the union of the source cell regions of these
full components, without filling any complementary pocket. On
$\mathcal T$, each chosen minus witness prevents the source on its
away side from reaching the target disk. Thus the entire $F_h^0$,
including its frontier, misses that disk. The transcript also
certifies the selected minus separations: the other endpoint of
each such pair is outside its explored complement-of-$W^-$
component. Every $N$ edge belongs to that complement graph, so these
certificates imply the corresponding $N$ disconnections. Denote them by
$\mathcal C(h)$.

In the residual law, let $\mathcal U_h$ be the event that there
exists an actual plus target loop enclosing exactly $I$, whose
closed disk misses $F_h^0$, together with actual plus witnesses
for the specified missing outside pairs of plus type. This event
is increasing in \eqref{eq:spatial-order}. The minimality rule and
the identities of the previously selected loops have been dropped;
they would not be increasing. If $\nu_g$ is the law of the
transcript under the original law conditioned on the core, the implication
from $\mathcal T$ gives
\[
 a_0\le\int \ind_{\mathcal C(h)}\,
                  \mu_h(\mathcal U_h)\,\nu_g(\dd h).
\]
It follows that the histories satisfying $\mathcal C(h)$ and
$\mu_h(\mathcal U_h)\ge a_0/2$ have $\nu_g$ mass at least $a_0/2$.
We work separately under each such $\mu_h$. There is no FKG or
circuit estimate under the mixture conditioned on this choice of
histories.

For this fixed history choose the accessible component $U_h$ of
the complement of the raw region $F_h^0$ containing the already
fixed mark $i_0$. Histories with $i_0\in F_h^0$ have
$\mu_h(\mathcal U_h)=0$ and do not occur in the retained family.
Every target disk in $\mathcal U_h$ is connected, contains $i_0$,
and misses $F_h^0$, so it lies in this one component $U_h$.
Only after choosing $U_h$, add the other complementary pockets
to the forbidden cell region and call the resulting region $F_h$.
This addition is for the marginal collar test alone. It does not
discard any actual residual variable from $\mu_h$; plus witnesses
for other pairs may still lie in those pockets.

For each individual connected source, the connected set $U_h$
lies in one component of its complement and hence faces one
incident boundary walk of that source. It may face that walk at
many entrances, but there are at most $k$ source labels in total.
Against $U_h$ the frontier is actual minus in the middle band.
The exact completion of one component of \cref{prop:circuits} therefore
applies there; any use of a partition comparison has the same
$k$ indivisible labels. The added pockets attach to these sources,
so every local piece in the band is anchored back to a source mark
in the old core or to the outer boundary. A target disk remains
disjoint from $F_h$.

Put $b_0=a_0/2$. Choose a collar width $\eta>0$, uniformly in the
retained histories, so that \cref{lem:collar} bounds by $b_0/2$
the existence of a full underlying dual traversal of the middle
band which lies in $F_h^\eta$ and avoids $F_h$ and its attached
minus circuits. Use the selected $U_h$ marginal for this estimate.
The full residual law $\mu_h$ is retained for all other events.
The virtual-disk clause makes this a simultaneous bound for
traversals inside any target disk in $\mathcal U_h$. In particular,
we subtract this error from $\mu_h(\mathcal U_h)$; we do not
condition the law on absence of the collar event.

We spell out the increasing event supplied by this subtraction.
Write $K$ for the enlarged lattice core, $A$ for its fixed middle
band, and $A^+$ for the wider band. Choose a fixed regular outer
dual graph $R$ inside $B$, containing $K$ with a complete connected
annular
shell outside $K$. In the path tests below every underlying dual
edge of $R$ is available, regardless of its $N$ state or whether
it crosses a minus edge or the forbidden region. For each required
pair across $I$, orient paths from its $I$ mark, and require every
simple underlying path $\gamma$ between the two marks to satisfy
\begin{equation}\label{eq:spatial-temporary-alternative}
 \begin{array}{c}
 \gamma\text{ crosses an actual }W^+\text{ edge in }K,\\
 \text{or}\\
 \gamma\text{ has a subpath }\beta\subset A
 \text{ from the inner to the outer boundary of }A\\
 \text{which visits a point }z\text{ with }\dist(z,F_h)>\eta .
 \end{array}
\end{equation}
Intersect this requirement with the actual plus separations for
the specified outside pairs, and call the event $\mathcal V_h$.
The path family is finite and fixed by $R$; $F_h$ is fixed by $h$.
Only the first alternative changes with $W^+$. Thus
$\mathcal V_h$ is increasing. Neither $F_h$ nor any minus edge is
a real barrier in this definition.

On $\mathcal U_h$ outside the collar event, \eqref{eq:spatial-temporary-alternative}
holds. Indeed a path avoiding the actual $W^+$ edges in $K$ must
exit $K$ before its first crossing of the target loop. An edge
meeting the closed core is counted in $K$, so a first crossing
on its boundary cannot be lost. The path prefix before that
target crossing lies in the target disk. Before its first exit
from $K$, take the segment from its last hit of the inner band
boundary to its first subsequent hit of the outer one. This is
a full subpath $\beta\subset A$, even if the original path made
earlier excursions. It lies in the target disk and therefore
avoids $F_h$ and all attached minus circuits. Absence of the
collar event forces it to visit a point at distance greater than
$\eta$. Hence
\begin{equation}\label{eq:spatial-temporary-probability}
                   \mu_h(\mathcal V_h)\ge b_0/2.
\end{equation}

We replace the second alternative by actual plus circuits. Choose
a deterministic grid in $A$ of mesh much smaller than $\eta$ and
than the separation of its two boundary curves. For each cell
whose filled buffer is disjoint from $F_h$ and contained in $A^+$,
request an actual plus circuit surrounding the cell in that buffer.
Choose the grid and buffers so that every point of $A$ at distance
greater than $\eta$ from $F_h$ lies in an inner cell surrounded by
one requested circuit, and the disk of every such circuit has diameter
smaller
than the distance between the two boundary curves of $A$. This
is possible by taking the cell size a sufficiently small fixed
multiple of $\eta$ and of the band width. There are at most
$M=M(\eta)$ cells. Their buffers miss all source queries and the
old core; for small mesh they exceed the lattice cutoff.

Under the full residual pin law $\mu_h$, each requested actual
plus circuit has probability at least $p_1>0$ by
\cref{prop:circuits}. Positive association intersects all these
increasing events with $\mathcal V_h$, giving probability at least
$(b_0/2)p_1^M$. On their intersection, no across-$I$ path can
avoid the actual plus edges in $K$. Otherwise
\eqref{eq:spatial-temporary-alternative} supplies a full band
subpath visiting the interior of one requested circuit. A connected
path avoiding that circuit and visiting its interior must stay
inside its disk. Its two band-boundary endpoints cannot both lie
there, by the diameter choice. This contradiction realizes the
temporary alternative using actual plus edges only.

For completeness, separation in the graph $R$ gives actual closed
witness loops in $K$. Any path using a more distant exterior can
have each exterior excursion replaced in the complete connected
shell of $R$ outside $K$. Thus separation in $R$ is separation in
the full underlying dual graph by $W^+\cap K$. Planar cut-cycle
duality yields a simple plus loop in $K$ for each across pair;
a boundary-to-boundary cut is impossible because these blocking
edges lie strictly inside $R$. Artificial fill and minus edges
play no role in this final separation.

The history certificates protect the outside minus pairs. The
outside plus pairs are protected by $\mathcal V_h$, the across
pairs by the new local plus loops, and there was no missing pair
wholly inside $I$. Previously short loops are fixed in $g$.
Therefore compatibility holds and at least one missing pair has
joined $\mathcal S(P_{j+1})$. The retained histories have
$\nu_g$ mass at least $b_0$, so the construction proves
\eqref{eq:spatial-witness-step} with, for example,
\[
             r(a,k)=\frac{b_0^2}{2}\,p_1^{M(\eta)}>0,
       \qquad b_0=\frac{a}{2L(k)}.
\]
The collar width, grid, and mesh threshold were chosen from
$a,k$ and the fixed bands, not from the particular core pattern.
Replace $r(a,k)$ by its minimum over the finitely many band groups.

We finally iterate without retaining the long source histories.
Let $\mu$ be the original law conditional on the initial pattern,
and set
\[
 \mathcal G_j=\{\text{compatibility and }|\mathcal S(P_j)|\ge j\}.
\]
Initially $\mu(\mathcal G_0)\ge p$. If
$\mu(\mathcal G_j)\ge q_j>0$, disintegrate \emph{this original law}
only by the complete pattern $P_j$. Patterns with
$|\mathcal S(P_j)|\ge j$ and posterior compatibility at least
$q_j/2$ have total mass at least $q_j/2$: the other patterns
contribute at most $q_j/2$ to $\mu(\mathcal G_j)$. If such a
pattern already has $j+1$ short pairs, its conditional contribution
is at least $q_j/2$; otherwise apply
\eqref{eq:spatial-witness-step}. Hence
\[
 \mu(\mathcal G_{j+1})\ge
       \frac{q_j}{2}\min\left\{\frac{q_j}{2},
                                  r(q_j/2,k)\right\}
             =:q_{j+1}>0.
\]
At the next step all source transcripts and selected type events
have been integrated out. The new conditional law is again the
original law given only its complete core pattern, so its unused
buffers have the pin/FKG law required above. After $m$ steps all
pairs have short witnesses in $K_m\Subset B$, with probability
at least $q_m$. All widths and the finitely many mesh thresholds
are chosen before taking the mesh limit. This proves the lemma.
\end{proof}

\begin{proposition}[Full-spin contiguity]
\label{prop:full-contiguity}
Fix compactly nested regular boxes $Q\Subset B$ in the interior
of a $W$-flat domain, with fixed positive buffer ratios. Thus the
complementary $N$ components are free. Let $\mu_\delta$
be its law and $\pi_\delta$ the plane law, restricted to the full
spin pair in $Q$. The domains and any allowed exterior $(s,W)$
data may depend on $\delta$, but their walls and prescriptions stay
outside $B$. There is a uniform contiguity modulus: for every
$\eps>0$ there is $\zeta>0$ such that, for every sufficiently small
mesh and every event $E$ determined by the full spin pair in $Q$,
\begin{equation}\label{eq:spatial-contiguity}
             \pi_\delta(E)<\zeta
           \quad\Longrightarrow\quad
             \mu_\delta(E)<\eps.
\end{equation}
The conclusion also holds if the tested data include $W$ in $Q$.
In particular it applies to mesh-dependent tests. Either flat
parity can be used by assigning $W$ to the parity carrying that
flat wall. This statement does not include a fixed-color wired
$N$ component before its removal.
\end{proposition}

\begin{proof}
Choose $Q\Subset Q'\Subset B'$ with
$B'\Subset B$, leaving the fixed buffers of
\cref{lem:buffered-witnesses}. Let $P$ be the complete $(s,W)$ pattern
in $Q'$, and let $T(P)$ be its number of terminals. By
\cref{prop:cluster-counts}, uniformly in both domain and plane laws,
\begin{equation}\label{eq:spatial-terminal-tail}
                       \P[T(P)>M]\le C e^{-cM}.
\end{equation}
Choose $M$ so that the right side is less than $\eps/4$ for the
domain law.

On $T(P)\le M$, write $b$ for an assignment of signs to its
terminals, and let $v_\mu(b\mid P)$ be its conditional exterior
compatibility probability. By \cref{lem:finite-cuts}, the actual
terminal-color probability $q_\mu(b\mid P)$ satisfies
\begin{equation}\label{eq:spatial-color-cost}
        2^{-M}v_\mu(b\mid P)
             \le q_\mu(b\mid P)\le v_\mu(b\mid P).
\end{equation}
Indeed, given a compatible exterior partition, specifying these
terminal signs costs $2^{-j}$ for some $j\le M$ distinct free global
components. The same calculation applies in the plane because
\cite[Theorem~17.1]{DKLM2026} gives no unbounded complementary
component and independent fair transverse colors. Components of
the restricted pattern not incident to an unprescribed edge carry
independent fair colors. Their distribution is the same under the
two laws once $P$ and the terminal signs are given. The lower
bound in \eqref{eq:spatial-color-cost} uses freeness; a terminal
wired to a prescribed incompatible color would not satisfy it.

Choose $p>0$ with $2^M p<\eps/4$. The total domain probability of
assignments with $v_\mu(b\mid P)<p$ is at most $2^M p$, by the
upper bound in \eqref{eq:spatial-color-cost}. For every remaining
assignment, \cref{lem:buffered-witnesses} gives short witnesses in
$B'$ with conditional probability at least $c(p,M)>0$.
Fix a realized pattern $g$ for $P$, and let $\mathcal W_b$ be this
short-witness event, determined by a
complete pattern in $B'$. If $C_{B'}$ and $C_{Q'}$ are the
unconditioned pattern density constants in
\cref{prop:pattern-comparison}, summing over the larger patterns
and then dividing gives
\[
 \pi_\delta(\mathcal W_b\mid P=g)
   =\frac{\pi_\delta(\mathcal W_b\cap\{P=g\})}{\pi_\delta(P=g)}
   \ge (C_{B'}C_{Q'})^{-1}
          \mu_\delta(\mathcal W_b\mid P=g).
\]
This uses only the original laws on one parity on the two pattern
sets, not a comparison after conditioning on $d$ or on the
witness event. Hence there is $c_0=c_0(p,M)>0$ such that
\begin{equation}\label{eq:spatial-plane-viability}
 v_\pi(b\mid P)\ge c_0,
 \qquad q_\pi(b\mid P)\ge 2^{-M}c_0
\end{equation}
whenever $v_\mu(b\mid P)\ge p$ and $T(P)\le M$.
There is no conditioning on $d$ in this comparison: the short
witness event concerns only $W$ and $s$.

Let $C_0$ be the uniform density-ratio constant for $P$ in
\cref{prop:pattern-comparison}. Conditional on $P$ and a terminal
assignment, integrate the identical independent internal colors.
On the retained assignments,
\eqref{eq:spatial-plane-viability} bounds the ratio of terminal-color
probabilities by $2^M/c_0$. Therefore every tested event satisfies
\begin{equation}\label{eq:spatial-contiguity-modulus}
       \mu_\delta(E)
       \le \frac{\eps}{2}
                + C_0\frac{2^M}{c_0}\,\pi_\delta(E).
\end{equation}
Choosing $\zeta<\eps c_0/(2C_0 2^M)$ proves
\eqref{eq:spatial-contiguity}. All constants were chosen from
$\eps$, the fixed geometry, and $t$ before the mesh limit, so the
modulus is uniform in the stated domain laws. Enlarging the tested
data by $W$ does not change the conditioning argument.
\end{proof}

The distinction between \cref{prop:pattern-comparison} and
\cref{prop:full-contiguity} is useful later. The former supplies
uniform multiplicative estimates even for extremely rare
patterns on one parity. The latter supplies domination in probability
for both spin parities, without asserting a bounded density ratio
for every possible microscopic color assignment.
